%% kapa: exam=1 round=35 year=2024 period=2 form=A subject=accounting subject_ko=회계학 seq=33 number=73
%% source: Q-Net 감정평가사 기출문제 (cst00309) artlSeq=5213229 file=raw/1008/5213229/2교시 문제지 원본.pdf
%% conversion: kordoc 4.24.51 -> tools/gen1.py (auto); review=proofread-sonnet; flags=tab
\begin{question}{73}
(주)감평은 종합원가계산제도를 채택하고 있으며, 제품 X의 생산관련 자료는 다음과 같다.
\par{\small\centering\begin{tabular}{|c|c|}\hline
구 분 & 물 량 \tabularnewline\hline
기초재공품(전환원가 완성도) & 60단위(70\%) \tabularnewline\hline
당기착수량 & 300단위 \tabularnewline\hline
기말재공품(전환원가 완성도) & 80단위(50\%) \tabularnewline\hline
\end{tabular}\par}
직접재료는 공정 초에 전량 투입되고, 전환원가(conversion cost, 또는 가공원가)는 공정 전반에 걸쳐 균등하게 발생한다. 품질검사는 전환원가(또는 가공원가) 완성도 80\% 시점에 이루어지며, 당기에 품질검사를 통과한 합격품의 5\%를 정상공손으로 간주한다. 당기에 착수하여 완성된 제품이 200단위일 때 비정상공손 수량은? (단, 재고자산의 평가방법은 선입선출법을 적용한다.)
\begin{choices}
  \item 7단위
  \item 10단위
  \item 13단위
  \item 17단위
  \item 20단위
\end{choices}
\end{question}
\answer{1}
